% Exercises for Chapter 13 (Gaussian vectors). Statements only: solutions are kept privately.
% Throughout, (Omega, F, P) is a probability space, vectors are columns, and N_N(mu, C) is the
% Gaussian law of the notes. The computational exercises can be done with probnotes.gaussian.

\begin{exercise}[Which matrices are covariance matrices?]\label{ex:gv-covariance-matrices}
\easy
\begin{enumerate}[label=(\alph*)]
  \item For $c\in\mathbb R$ let
    $C_c=\begin{pmatrix}2&1&-1\\1&2&1\\-1&1&c\end{pmatrix}$. Show that $C_c$ is the covariance
    matrix of some random vector if and only if $c\ge2$. For $c=1$ (the matrix written in the
    handwritten version of \noteref{exa:gv-affine}), find $\underline a\in\mathbb R^3$ with
    $\underline a^TC_1\underline a<0$, and explain directly, without
    \noteref{lem:gv-covariance-psd}, why no random vector $\underline X$ with components in $L^2$
    has covariance matrix $C_1$.
  \item Let $\underline X\sim\mathcal N_3(\underline 0,C_2)$. Find $\underline a\ne\underline 0$ with
    $C_2\underline a=\underline 0$, and show that $\underline a^T\underline X=0$ almost surely. On
    which plane does the law of $\underline X$ live?
  \item Let $N\ge2$, $\rho\in\mathbb R$, and $R_\rho\in\mathcal M_N(\mathbb R)$ the matrix with $1$
    on the diagonal and $\rho$ everywhere else. Show that $R_\rho$ is a covariance matrix if and only
    if $-\frac1{N-1}\le\rho\le1$. \emph{Hint:} find the eigenvalues of $R_\rho$, using the vector
    $(1,\dots,1)^T$ and the vectors orthogonal to it.
  \item For $0\le\rho\le1$, let $W_0,W_1,\dots,W_N$ be \iid with law $\mathcal N(0,1)$ and
    $X_i=\sqrt\rho\,W_0+\sqrt{1-\rho}\,W_i$. Show that $\underline X\sim\mathcal N_N(\underline 0,R_\rho)$,
    giving $A\in\mathcal M_{N\times(N+1)}(\mathbb R)$ with $AA^T=R_\rho$. For
    $\rho=-\frac1{N-1}$, show that $X_1+\dots+X_N=0$ almost surely for every
    $\underline X\sim\mathcal N_N(\underline 0,R_\rho)$.
\end{enumerate}
\end{exercise}

\begin{exercise}[Gaussian vectors through linear combinations]\label{ex:gv-linear-combinations}
\medium
\begin{enumerate}[label=(\alph*)]
  \item Let $A\in\mathcal M_{N\times K}(\mathbb R)$ and $A'\in\mathcal M_{N\times K'}(\mathbb R)$ with
    $AA^T=A'A'^T$, and let $\underline X\sim\mathcal N_K(\underline 0,I_K)$,
    $\underline X'\sim\mathcal N_{K'}(\underline 0,I_{K'})$. Show that $A\underline X+\underline\mu$
    and $A'\underline X'+\underline\mu$ have the same law. Find $A'\in\mathcal M_{2\times1}(\mathbb R)$
    with $A'A'^T=AA^T$ for the matrix $A=\left(\begin{smallmatrix}1&1\\2&2\end{smallmatrix}\right)$ of
    \noteref{exa:gv-degenerate}, and describe the two maps $\underline x\mapsto A\underline x$ and
    $x\mapsto A'x$.
  \item Let $\underline Y$ be a random vector with values in $\mathbb R^N$. Show that
    $\underline Y$ is a Gaussian vector if and only if, for every $\underline a\in\mathbb R^N$, the
    real random variable $\underline a^T\underline Y$ is Gaussian (possibly degenerate).
    \emph{Hint:} for ``$\Leftarrow$'', show first that the $Y_i$ are in $L^2$, call
    $\underline\mu$ and $C$ the mean vector and the covariance matrix of $\underline Y$, and compute
    $\Phi_{\underline Y}(\underline a)=\Phi_{\underline a^T\underline Y}(1)$.
  \item Let $X_1,\dots,X_N$ be independent Gaussian random variables, some of which may be
    degenerate. Show in two ways that $(X_1,\dots,X_N)^T$ is a Gaussian vector: with (b), and with
    \noteref{prop:gv-concatenation}.
\end{enumerate}
\end{exercise}

\begin{exercise}[Degenerate Gaussian vectors]\label{ex:gv-degenerate}
\medium This exercise proves \noteref{rem:gv-degenerate}, which the notes state without proof. Let
$\underline Y=A\underline X+\underline\mu$ with $\underline X\sim\mathcal N_K(\underline 0,I_K)$,
$A\in\mathcal M_{N\times K}(\mathbb R)$, $C=AA^T$, and let $r=\operatorname{rank}C$.
\begin{enumerate}[label=(\alph*)]
  \item Show that $\ker C=\ker A^T$, and deduce that $\operatorname{Im}C=\operatorname{Im}A$.
  \item Show that $\mathbb P(\underline Y\in\underline\mu+\operatorname{Im}C)=1$.
  \item Suppose that $C$ is singular. Show that a proper linear subspace $V$ of $\mathbb R^N$ has
    Lebesgue measure zero. (\emph{Hint:} $V$ is contained in a hyperplane
    $\{\underline x:\underline a^T\underline x=0\}$ with $|\underline a|=1$; complete $\underline a$
    to an orthonormal basis and change variables with \noteref{thm:rv-jacobi} in $\mathbb R^N$.)
    Deduce that $\underline Y$ is not absolutely continuous, and conclude, with
    \noteref{thm:gv-density}, that a Gaussian vector is absolutely continuous if and only if its
    covariance matrix is invertible.
  \item Let $r\ge1$, and let $U\in\mathcal M_{N\times r}(\mathbb R)$ have orthonormal columns that
    span $\operatorname{Im}C$, so that $C=U\Lambda U^T$ with $\Lambda$ diagonal and positive definite
    (spectral theorem). Show that $\underline W=U^T(\underline Y-\underline\mu)\sim\mathcal N_r(\underline 0,\Lambda)$
    is absolutely continuous on $\mathbb R^r$, and that $\underline Y=\underline\mu+U\underline W$
    almost surely: a degenerate Gaussian vector is a non-degenerate one of dimension $r$, placed in
    $\mathbb R^N$ by an affine map. Find $U$, $\Lambda$ and the law of $\underline W$ for
    \noteref{exa:gv-degenerate}.
\end{enumerate}
\end{exercise}

\begin{exercise}[An atom at zero]\label{ex:gv-random-sign}
\easy This is a computational exercise on the pair $(X,Y)=(X,ZX)$ of \noteref{exa:gv-not-joint}.
\begin{enumerate}[label=(\alph*)]
  \item Show that $\operatorname{cov}(X,Y)=0$, but that $X$ and $Y$ are not independent. Why does
    this not contradict \noteref{thm:gv-independence}?
  \item Show that the law of $X+Y$ is $\frac12\delta_0+\frac12\mathcal N(0,4)$, and write its
    distribution function $F$. Where is $F$ discontinuous?
  \item Simulate $n=10^5$ pairs (seed $2026$). Give an estimate of $\mathbb P(X+Y=0)$ with an
    asymptotic $95\%$ confidence interval. Compute $\sup_x|F_n(x)-F(x)|$ for the empirical
    distribution function $F_n$ of the simulated values of $X+Y$ (careful: there are many ties at
    $0$; compare $F_n$ and $F$ at each sample point, from the right and from the left), and compare
    it with the half-width $\varepsilon=\sqrt{\log(2/0.05)/(2n)}$ of the $95\%$
    Dvoretzky--Kiefer--Wolfowitz band. Do the same for $Y$ against the $\mathcal N(0,1)$ distribution
    function.
  \item Compute the sample correlation of $X$ and $Y$, and of $X^2$ and $Y^2$. What do they say?
\end{enumerate}
\end{exercise}

\begin{exercise}[Uncorrelated, Gaussian, and still dependent]\label{ex:gv-sign-flip}
\medium Let $X\sim\mathcal N(0,1)$, $c>0$, and $Y=X$ on $\{|X|\le c\}$, $Y=-X$ on $\{|X|>c\}$.
Write $\varphi$ and $\Phi$ for the density and the distribution function of $\mathcal N(0,1)$.
\begin{enumerate}[label=(\alph*)]
  \item Show that $Y\sim\mathcal N(0,1)$.
  \item Show that $\operatorname{cov}(X,Y)=2\big(2\Phi(c)-1-2c\varphi(c)\big)-1$, that this is an
    increasing function of $c$ with limits $-1$ and $1$, and find numerically the unique $c^*$ with
    $\operatorname{cov}(X,Y)=0$.
  \item For $c=c^*$, show that $X$ and $Y$ are not independent, and deduce that $(X,Y)^T$ is not a
    Gaussian vector. Check this directly: compute $\mathbb P(X+Y=0)$.
  \item Simulate $10^6$ pairs with $c=c^*$ (seed $2026$). Is the sample correlation within two
    standard errors of $0$? (\emph{Careful:} what is the standard error of the sample mean of $XY$
    here?) Compare the fraction of pairs with $X+Y=0$ with (c).
\end{enumerate}
\end{exercise}

\begin{exercise}[The bivariate normal law]\label{ex:gv-bivariate}
\medium (in the style of Ross, Ch.~6) Let $\mu_X,\mu_Y\in\mathbb R$, $\sigma_X,\sigma_Y>0$,
$-1<\rho<1$, and
\[
  (X,Y)^T\sim\mathcal N_2\left(\begin{pmatrix}\mu_X\\\mu_Y\end{pmatrix},
  \begin{pmatrix}\sigma_X^2&\rho\sigma_X\sigma_Y\\\rho\sigma_X\sigma_Y&\sigma_Y^2\end{pmatrix}\right).
\]
\begin{enumerate}[label=(\alph*)]
  \item Show that the covariance matrix is invertible, that $\rho$ is the correlation coefficient
    of $X$ and $Y$, and that
    \[
      f_{(X,Y)}(x,y)=\frac1{2\pi\sigma_X\sigma_Y\sqrt{1-\rho^2}}
      \exp\Big[-\frac{u^2-2\rho uv+v^2}{2(1-\rho^2)}\Big],
    \]
    where $u=(x-\mu_X)/\sigma_X$ and $v=(y-\mu_Y)/\sigma_Y$.
    When are $X$ and $Y$ independent?
  \item The scores of a student in two exams are modelled by such a pair, with $\mu_X=70$,
    $\sigma_X=10$, $\mu_Y=65$, $\sigma_Y=12$ and $\rho=0.6$. Find the law of the average score
    $\frac12(X+Y)$, the probability that the total $X+Y$ exceeds $150$, and the probability that the
    student does better in the second exam.
  \item In (b), find $a\in\mathbb R$ such that $X+aY$ and $X-Y$ are independent.
  \item Show that, for $r>0$, the ellipse
    $E_r=\{(x,y):u^2-2\rho uv+v^2\le(1-\rho^2)r^2\}$ has probability
    $\mathbb P\big((X,Y)\in E_r\big)=1-e^{-r^2/2}$, whatever the parameters. \emph{Hint:} whiten as
    in the proof of \noteref{thm:gv-density}, and use \noteref{prop:rv-notable-sums}~d). Which $r$
    gives an ellipse of probability $0.9$?
\end{enumerate}
\end{exercise}

\begin{exercise}[Sample mean and sample variance]\label{ex:gv-sample-mean}
\hard (in the style of Ross, Ch.~7) Let $X_1,\dots,X_n$ be \iid with law $\mathcal N(\mu,\sigma^2)$,
$\sigma>0$, $n\ge2$, and let
\[
  \bar X=\frac1n\sum_{i=1}^nX_i,\qquad S^2=\frac1{n-1}\sum_{i=1}^n(X_i-\bar X)^2 .
\]
\begin{enumerate}[label=(\alph*)]
  \item Show that $(\bar X,X_1-\bar X,\dots,X_n-\bar X)^T$ is a Gaussian vector of
    $\mathbb R^{n+1}$, and compute its mean vector and its covariance matrix. Is it absolutely
    continuous?
  \item Deduce that $\bar X$ is independent of $(X_1-\bar X,\dots,X_n-\bar X)$, and hence of $S^2$.
  \item Let $Q\in\mathcal M_n(\mathbb R)$ be an orthogonal matrix whose first row is
    $\frac1{\sqrt n}(1,\dots,1)$, and $\underline W=Q\underline X$. Find the law of $\underline W$, and
    show that $\sum_{i=1}^n(X_i-\bar X)^2=\sum_{k=2}^nW_k^2$. Deduce that
    $(n-1)S^2/\sigma^2\sim\chi^2(n-1)$, and compute $\mathbb E[S^2]$ and $\operatorname{var}(S^2)$.
  \item The independence in (b) is special to Gaussian samples. Let $n=2$ and $X_1,X_2$ be \iid
    with law $\mathrm{Exp}(1)$. Show that $\operatorname{cov}(\bar X,S^2)=1$.
\end{enumerate}
\end{exercise}

\begin{exercise}[Sampling Gaussian vectors, also degenerate ones]\label{ex:gv-sampling}
\medium This is a computational exercise. Let $C$ be the covariance matrix of
\noteref{exa:gv-affine} and $C_2$ the singular matrix of \exref{ex:gv-covariance-matrices}.
\begin{enumerate}[label=(\alph*)]
  \item Simulate $n=10^5$ points of $\mathcal N_3(\underline 0,C)$ (seed $2026$) twice, as
    $A\underline X$ with $A$ the symmetric square root of $C$ and with $A$ its Cholesky factor. For
    each sample, compare the sample covariance matrix with $C$, and the sample covariance matrix of
    $\underline Z=(X_1-X_2,X_1+X_3)^T$ with the matrix $\Sigma$ of \noteref{exa:gv-affine}.
  \item Show that if $\underline Y\sim\mathcal N_N(\underline\mu,C)$ with $C$ invertible, then
    $(\underline Y-\underline\mu)^TC^{-1}(\underline Y-\underline\mu)\sim\chi^2(N)$. For the samples of
    (a), compare the fraction of points with $\underline y^TC^{-1}\underline y\le r$ with the
    $\chi^2(3)$ distribution function at $r=1$, $3$ and $7.815$.
  \item Check that \texttt{numpy.linalg.cholesky} fails on $C_2$. Sample $10^5$ points of
    $\mathcal N_3(\underline 0,C_2)$ with the symmetric square root and with a Cholesky
    factorisation that allows zero pivots (\texttt{probnotes.gaussian.psd\_cholesky}, or your own),
    and check that all the points lie on the plane of \exref{ex:gv-covariance-matrices}~(b).
  \item How large must $n$ be for the sample variance of the first coordinate (whose variance is
    $2$) to have standard deviation at most $0.01$? Use \exref{ex:gv-sample-mean}~(c). Check your
    answer by computing $1000$ independent sample variances of that size (seed $2026$).
\end{enumerate}
\end{exercise}

\begin{exercise}[Whitening a Gaussian sample]\label{ex:gv-whitening}
\medium This is a computational exercise. Let $\underline Y_1,\dots,\underline Y_n$ be \iid with law
$\mathcal N_N(\underline\mu,C)$, $C$ invertible.
\begin{enumerate}[label=(\alph*)]
  \item Show that if $A\in\mathcal M_N(\mathbb R)$ and $AA^T=C$, then the vectors
    $A^{-1}(\underline Y_k-\underline\mu)$ are \iid with law $\mathcal N_N(\underline 0,I_N)$. Show
    that if also $B\in\mathcal M_N(\mathbb R)$ and $BB^T=C$, then $B=AQ$ for an orthogonal matrix $Q$: two ``whitenings'' differ by an
    orthogonal transformation.
  \item In practice $\underline\mu$ and $C$ are unknown, and one uses the sample mean
    $\bar{\underline Y}$ and the sample covariance matrix $\hat C$ (with divisor $n-1$), assumed
    invertible. Show that the vectors $\hat C^{-1/2}(\underline Y_k-\bar{\underline Y})$, with
    $\hat C^{-1/2}$ the inverse of the symmetric square root of $\hat C$, have sample mean exactly
    $\underline 0$ and sample covariance matrix exactly $I_N$, whatever the law of the
    $\underline Y_k$.
  \item Simulate $n=2000$ points with $\underline\mu=(1,-2,0.5)^T$ and $C$ the matrix of
    \noteref{exa:gv-affine} (seed $2026$). Whiten them with the true parameters, by the inverse
    symmetric square root and by the inverse Cholesky factor, and with the estimated ones. Report the
    largest off-diagonal entry of the sample covariance matrix in each case, and find the orthogonal
    matrix that relates the two whitenings with the true parameters.
  \item With the true parameters, each off-diagonal entry of the sample covariance matrix of the
    whitened sample has standard deviation close to $n^{-1/2}$: why? How large must $n$ be for the
    $95\%$ asymptotic interval $\pm1.96\,n^{-1/2}$ to be at most $\pm0.01$? Check by simulation how
    often the entry falls within $0.01$, in $100$ repetitions (one generator, seed $2026$) with
    samples of that size from $\mathcal N_2\big((1,-2)^T,\left(\begin{smallmatrix}2&1\\1&2\end{smallmatrix}\right)\big)$.
\end{enumerate}
\end{exercise}
